% =====================================================================
%  Section 02 — Phase 1 introduction
% =====================================================================
\makesectioncover{Phase 1 \textbullet\ The Existing Reality}

\section{Phase 1 \textbullet\ Auditing the Existing System}
\label{sec:phase1-intro}

\subsection{What Phase 1 is doing and why}

Phase 1 is the audit phase. Before we can ask whether new infrastructure
is solving real problems we need to know what the existing transport
reality \emph{is}.

The work splits across two notebooks:

\begin{itemize}
  \item \fn{Cleaning Data.ipynb} — loads seven raw GeoJSON datasets,
        reprojects them to a metric CRS, KNN-imputes missing values,
        and builds seven spatial/attribute merges (A--G) that downstream
        analysis depends on.
  \item \fn{Visual Analysis.ipynb} — answers twelve research questions
        (Q1--Q12) using those merges and produces 40 Plotly visualisations.
        The questions converge on \kw{four structural gaps} that frame
        every Phase 2 hypothesis: ghost terminals, empty-return terminals,
        vehicle-route mismatch, underserved hexagons.
\end{itemize}

\subsection{The two-system framing}

Throughout this report we use two labels:

\begin{itemize}
  \item \kw{System A · formal} — the published, official network. In Phase 1
        this is approximated by the routes-with-fares (\fn{Public Transport
        Routes.json}) and the published terminals
        (\fn{Public Transport Terminals.json}). In Phase 2 it is augmented
        with GTFS, the metro network, the LRT, and the BRT.
  \item \kw{System B · informal/lived} — the network the city actually
        rides every morning. In Phase 1 this is approximated by the
        microbus boardings (\fn{Cairo Daily Boarding.json}) and the
        passenger-flow segments (\fn{Public Transport Passenger Flow.json}).
\end{itemize}

The point of Phase 1 is to make the two systems comparable on a single
map and a single set of metrics — buffer distances, boarding totals,
fare-per-km. The four structural gaps are the cells of that comparison
that fail.
